\section{Related Works}


% \subsection{Program Analysis for Smart Contracts}

% \noindent \textbf{Parametric Optimization.}
% For some flash-loan attack cases, prior work manually extracted formulas of candidate functions,
% manually defined parameter constraints, and applied off-the-shelf optimization to search for profitable parameters\cite{cao2021flashot,qin2021attacking}.
% These methods demonstrate feasibility for specific incidents, but their manual workflow and protocol-specific expertise requirements
% limit scalability when the attack space is large.

\noindent \textbf{Static Analysis and Symbolic Execution.}
Static analyzers such as Slither, Securify, Zeus, Park, and 
SmartCheck\cite{feist2019slither,tsankov2018securify,
kalra2018zeus,zheng2022park,tikhomirov2018smartcheck}
have been widely used to detect vulnerabilities in smart contracts.
Symbolic-execution-based tools and 
techniques\cite{king1976symbolic,mythril,luu2016making,liu2020towards,
frank2020ethbmc,feng2019precise,mossberg2019manticore}
systematically explore path conditions to find property-violating traces.
These approaches are highly effective for many vulnerability classes. 
However, for complex DeFi smart contracts, real-world attacks typically 
involves multiple contracts and long execution traces, 
which can lead to state coupling and path explosion 
challenges for static and symbolic methods.

\noindent \textbf{Fuzzing.}
ContractFuzzer, sFuzz, ContraMaster, SMARTIAN, 
and ItyFuzz\cite{jiang2018contractfuzzer,nguyen2020sfuzz,wang2020oracle,
choi2021smartian,shou2023ityfuzz}
introduce effective fuzzing strategies for smart-contract security testing.
However, these techniques either only work on one contract or focus
on specific vulnerabilities like re-entrancy. 
Moreover, real-world DeFi smart contract attacks can 
involve tens of contract interactions and tens of chosen parameters,
which is unlikely to be found by random fuzzing. 


\noindent \textbf{Smart Contract Attack Monitoring.}
Prior studies analyze DeFi smart contract attacks and related
 DeFi transaction strategies\cite{qin2021attacking,cao2021flashot,
 werapun2022flash,wang2021towards,gan2022understanding}.
Monitoring-oriented systems have also been proposed 
for post-hoc or real-time attack detection
\cite{xia2023detecting,ramezany2023midnight,wang2022defiscanner}.
Despite these efforts and the good coverage of known attack patterns, 
these works do not propose systematic generalized methods for defending 
against novel attack vectors that may arise in the future. 



\subsection{Runtime Enforcement and Invariant Guards for Smart Contracts}

\noindent \textbf{Invariant Generation and Enforcement.}
Prior research has proposed multiple approaches to 
derive and enforce smart-contract invariants.
Cider\cite{liu2022learning} learns arithmetic-overflow invariants 
using reinforcement learning,
InvCon and InvCon+\cite{liu2022invcon,liu2024automated} combine 
dynamic inference with static validation. 
These directions provide practical invariant families 
and evidence for runtime protection,
especially when the objective is to reject abnormal executions 
with low deployment friction. However, 
these works only cover a limited set of invariant types and 
more importantly, do not propose methods for 
patching existing smart contracts after their deployment, 
making their approach limited to pre-deployment use cases. 
Our thesis contributes a broader framework for runtime safety enforcement,
including pre-deployment and post-deployment use cases.


\noindent \textbf{Control-Flow Restriction in industry.}
Industry systems such as SphereX use control-flow 
restriction to restrict control flows over deployed smart contracts\cite{spherex}.
However, their approach requires 
customers to manually define the allowed control flows, 
which is not scalable for complex DeFi protocols. Our thesis 
contributes an automated approach 
to deploy template-based control-flow guards that can be applied 
to a wide range of DeFi protocols with minimal manual effort and minimal 
number of falsely blocked transactions.



\noindent
\textbf{Re-entrancy Attack Defense and Secure Type System.}
Numerous tools restrict control flows to combat re-entrancy attacks. 
Static analysis tools~\cite{luu2016making,torres2018osiris,
liu2018reguard,feist2019slither,
zhang2019mpro,bose2022sailfish,ma2021pluto,zheng2022park,wang2024efficiently,
albert2020taming} identify re-entrancy vulnerabilities and apply guards. 
Runtime frameworks like Sereum~\cite{rodler2018sereum} and 
\citeauthor{grossman2017online} protect deployed contracts, 
while \citeauthor{callens2024temporarily} prevent duplicate 
function calls within transactions. Unlike these 
re-entrancy-focused approaches, our thesis adopts broader 
control flow restriction targeting comprehensive vulnerability classes.
\citeauthor{cecchetti2020securing} propose a security type system to 
enforce information flow and re-entrancy controls~\cite{cecchetti2020securing,
cecchetti2021compositional}. Conversely, 
our thesis contributes a purely runtime system built entirely on the EVM, 
operating without supplementary type systems or language modifications.




